% chapterintro
% ==============================
% 第5章 导言：chapter5-0-chinese.md
STL不仅包含数据结构，还有很多算法。数据结构可以帮助存放特定情况下需要保存的数据，而算法则会将数据结构中存储的数据进行变换。

让我们来看一个标准的例子，例如对\texttt{vector}实例中的数据进行累加。这个可以简单的通过循环迭代\texttt{vector}中的元素，将所有值累加在一个对应的值上：

\begin{codebox}[cpp]
vector<int> v {100, 400, 200 /*, ... */ };

int sum {0};
for (int i : v) { sum += i; }

cout << sum << '\n';
\end{codebox}

不过，作为一个标准的例子，当然可以使用STL的算法来完成：

\begin{codebox}[cpp]
cout << accumulate(begin(v), end(v), 0) << '\n';
\end{codebox}

例子中循环变量也不是很长，不过其可读性比\texttt{accumulate}差很多。一个10行的循环代码看起来的确很尴尬，那么本章我们就看来了解一下标准算法(\texttt{accumulate}, \texttt{copy}, \texttt{move}, \texttt{transform}和\texttt{shuffle}等等)的工作机制。

其思想就是提供丰富的算法供开发者使用，避免耗费大量的时间在重复制造轮子上面。另一方面就是，即便开发者会自己去实现相应STL中的算法，也要进行大量的测试来确保自己实现的算法是否正确，STL提供的算法都是经过了严格的测试。所以没有必要做重复的工作，这样也能节省代码审阅者的时间，否则他们还要确定算法实现中是否有Bug。

另一个重点是STL算法非常的高效。很多STL算法提供了多种特化实现，这样足以应对依赖迭代器类型的使用方式。例如，将\texttt{vector}中的所有元素都填充0时，就可以使用\texttt{std::fill}。因为\texttt{vector}使用的是一段连续的内存，对于这种使用连续内存存放数据的结构都可以使用\texttt{std::fill}进行填充，这个函数类似于C中的\texttt{memset}函数。当开发者将容器类型从\texttt{vector}改为\texttt{list}，STL算法就不能再使用\texttt{memset}了，并且需要逐个迭代list的元素，并将元素赋0。开发者不能为了使用\texttt{memset}将数据类型写死为\texttt{vector}或\texttt{std::fill}0，因为实际项目中，还是有很多数据结构存储的地址并不是连续的。大多数情况下，想要自己去将代码实现的更聪明是没有太多意义的，因为STL的实现者已经考虑到了这种情况，并且STL还是免费使用的，为什么不用呢？

让我们总结一下前面提到的几点。使用STL算法的好处：

\begin{itemize}
\item \textbf{维护性}：算法的名字已经说明它要做什么了。显式使用循环的方式与使用STL算法的方式没法对比。

\item \textbf{正确性}：STL是由专家编写和审阅过的，并且经过了良好的测试，重新实现的复杂程度可能是你无法想象的。

\item \textbf{高效性}：STL算法真的很高效，至少要比手写的循环要强许多。

\end{itemize}

很多算法都是对迭代器进行操作，第3章已经解释了迭代器的工作原理。本章专注于如何使用STL算法解决各种问题，了解这些STL应该如何使用。要展示所有STL算法的使用方式不是本书所要做的事情，这个事情C++手册已经完成了，你可以在网上进行查询，或者花钱购买电子/纸质发布版本。

作为一个STL“忍者”需要将C++手册放在手边……嗯，至少放在浏览器的书签中吧。当我们在完成一个任务的过程中，每个开发者都可以回看一下任务本身，在完成自己的任务时，确定这个STL算法是否适合于你的问题。

在线版本的C++手册：\texthref{http://cppreference.com}

其也提供离线下载功能。

\begin{myInfoBox}
Note: 在面试过程中，对于STL算法的熟悉程度也是判断一个开发者对C++的熟悉程度的标准之一。
\end{myInfoBox}


% --- recipe 1: chapter5-1-chinese.md ---
\section{容器间相互复制元素}

大多数STL数据结构都支持迭代器。这就意味着大多数数据结构能够通过成员函数\texttt{begin()}和\texttt{end()}成员函数得到相应的迭代器，并能对数据进行迭代。迭代的过程看起来是相同的，无论是什么样的数据结构都是一样的。

我们可以对\texttt{vector}，\texttt{list}，\texttt{deque}，\texttt{map}等等数据结构进行迭代。我们甚至可以使用迭代器作为文件/标准输入输出的出入口。此外，如之前章节介绍，我们能将迭代器接口放入算法中。这样的话，我们可以使用迭代器访问任何元素，并且可以将迭代器作为STL算法的参数传入，对特定范围内的数据进行处理。

\texttt{std::copy}算法可以很好的展示迭代器是如何将不同的数据结构进行抽象，而后将一个容器的数据拷贝到另一个容器。类似这样的算法就与数据结构的类型完全没有关系了。为了证明这点，我们会把玩一下\texttt{std::copy}。

\mySmallsection{How to do it...}

本节中，我们将对不同的变量使用\texttt{std::copy}。

\begin{enumerate}
\item 首先，包含必要的头文件，并声明所用到的命名空间。

\begin{codebox}[cpp]
#include <iostream>
#include <vector>
#include <map>
#include <string>
#include <tuple>
#include <iterator>
#include <algorithm>
using namespace std;
\end{codebox}

\item 我们将使用整型和字符串值进行组对。为了能很好的将其进行打印，我们将会重载\texttt{<<}流操作：

\begin{codebox}[cpp]
namespace std {
ostream& operator<<(ostream &os, const pair<int, string> &p)
{
    return os << "(" << p.first << ", " << p.second << ")";
}
}
\end{codebox}

\item 主函数中，我们将使用整型-字符串对填充一个\texttt{vector}。并且我们声明一个\texttt{map}变量，其用来关联整型值和字符串值：

\begin{codebox}[cpp]
int main()
{
    vector<pair<int, string>> v {
        {1, "one"}, {2, "two"}, {3, "three"},
        {4, "four"}, {5, "five"}};
    map<int, string> m;
\end{codebox}

\item 现在将\texttt{vector}中的前几个整型字符串对使用\texttt{std::copy\_n}拷贝到\texttt{map}中。因为\texttt{vector}和\texttt{map}是两种完全不同的结构体，我们需要对\texttt{vector}中的数据进行变换，这里就要使用到\texttt{insert\_iterator}适配器。\texttt{std::inserter}函数为我们提供了一个适配器。在算法中使用类似\texttt{std::copy\_n}的算法时，需要与插入迭代器相结合，这是一种更加通用拷贝/插入元素的方式(从一种数据结构到另一种数据结构)，但这种方式不是最快的。使用指定数据结构的成员函数插入元素无疑是更加高效的方式：

\begin{codebox}[cpp]
    copy_n(begin(v), 3, inserter(m, begin(m)));
\end{codebox}

\item 让我们打印一下\texttt{map}中的内容。纵观本书，我们会经常使用\texttt{std::copy}函数来打印容器的内容。\texttt{std::ostream\_iterator}在这里很有用，因为其可以将用户的标准输出作为另一个容器，而后将要输出的内容拷贝过去：

\begin{codebox}[cpp]
    auto shell_it (ostream_iterator<pair<int, string>>{cout,
    ", "});
    copy(begin(m), end(m), shell_it);
    cout << '\n';
\end{codebox}

\item 对\texttt{map}进行清理，然后进行下一步的实验。这次，我们会将\texttt{vector}的元素\emph{移动}到\texttt{map}中，并且是所有元素：

\begin{codebox}[cpp]
    m.clear();
    move(begin(v), end(v), inserter(m, begin(m)));
\end{codebox}

\item 我们将再次打印\texttt{map}中的内容。此外，\texttt{std::move}是一种改变数据源的算法，这次我们也会打印\texttt{vector}。这样，我们就会看到算法时如何对数据源进行的移动：

\begin{codebox}[cpp]
    copy(begin(m), end(m), shell_it);
    cout << '\n';
    copy(begin(v), end(v), shell_it);
    cout << '\n';
}
\end{codebox}

\item 编译运行这个程序，看看会发生什么。第一二行非常简单，其反应的就是\texttt{copy\_n}和\texttt{move}算法执行过后的结果。第三行比较有趣，因为移动算法将其源搬移到\texttt{map}中，所以这时的\texttt{vector}是空的。在重新分配空间前，我们通常不应该访问成为移动源的项。但是为了这个实验，我们忽略它：

\begin{codebox}[cpp]
$ ./copying_items
(1, one), (2, two), (3, three),
(1, one), (2, two), (3, three), (4, four), (5, five),
(1, ), (2, ), (3, ), (4, ), (5, ),
\end{codebox}

\end{enumerate}

\mySmallsection{How it works...}

\texttt{std::copy}是STL中最简单的算法之一，其实现也非常短。我们可以看一下等价实现：

\begin{codebox}[cpp]
template <typename InputIterator, typename OutputIterator>
OutputIterator copy(InputIterator it, InputIterator end_it,
OutputIterator out_it)
{
    for (; it != end_it; ++it, ++out_it) {
        *out_it = *it;
    }
    return out_it;
}
\end{codebox}

这段代码很朴素，使用\texttt{for}循环将一个容器中的元素一个个的拷贝到另一个容器中。此时，有人就可能会发问："使用\texttt{for}循环的实现非常简单，并且还不用返回值。为什么要在标准库实现这样的算法？"，这是个不错的问题。

\texttt{std::copy}并非能让代码大幅度减少的一个实现，很多其他的算法实现其实非常复杂。这种实现其实在代码层面并不明显，但STL算法更多的在于做了很多底层优化，编译器会选择最优的方式执行算法，这些底层的东西目前还不需要去了解。

STL算法也让能避免让开发者在代码的可读性和优化性上做权衡。

\begin{myInfoBox}
Note： 如果类型只有一个或多个(使用\texttt{class}或\texttt{struct}包装)的矢量类型或是类，那么其拷贝赋值通常是轻量的，所以可以使用\texttt{memcopy}或\texttt{memmove}进行赋值操作，而不要使用自定义的赋值操作符进行操作。
\end{myInfoBox}

这里，我们也使用了\texttt{std::move}。其和\texttt{std::copy}一样优秀，不过\texttt{std::move(*it)}会将循环中的源迭代器，从局部值(左值)转换为引用值(右值)。这个函数就会告诉编译器，直接进行移动赋值操作来代替拷贝赋值操作。对于大多数复杂的对象，这会让程序的性能更好，但会破坏原始对象。


% --- recipe 2: chapter5-2-chinese.md ---
\section{容器元素排序}

排序是一项很常见的任务，并且可以通过各种各样的方式进行。每个计算机科学专业的学生，都学过很多排序算法(包括这些算法的性能和稳定性)。

因为这是个已解决的问题，所以开发者没必要浪费时间再次来解决排序问题，除非是出于学习的目的。

\mySmallsection{How to do it...}

本节中，我们将展示如何使用\texttt{std::sort}和\texttt{std::partial\_sort}。

\begin{enumerate}
\item 首先，包含必要的头文件和声明所使用的命名空间。

\begin{codebox}[cpp]
#include <iostream>
#include <algorithm>
#include <vector>
#include <iterator>
#include <random>
using namespace std;
\end{codebox}

\item 我们将打印整数在\texttt{vector}出现的次数，为了缩短任务代码的长度，我们在这里写一个辅助函数：

\begin{codebox}[cpp]
static void print(const vector<int> &v)
{
    copy(begin(v), end(v), ostream_iterator<int>{cout, ", "});
    cout << '\n';
}
\end{codebox}

\item 我们开始实例化一个\texttt{vector}：

\begin{codebox}[cpp]
int main()
{
    vector<int> v {1, 2, 3, 4, 5, 6, 7, 8, 9, 10};
\end{codebox}

\item 因为我们将使用不同的排序函数将\texttt{vector}多次打乱，所以我们需要一个随机数生成器：

\begin{codebox}[cpp]
    random_device rd;
    mt19937 g {rd()};
\end{codebox}

\item \texttt{std::is\_sorted}函数会告诉我们，容器内部的值是否已经经过排序。所以这行将打印到屏幕上：

\begin{codebox}[cpp]
    cout << is_sorted(begin(v), end(v)) << '\n';
\end{codebox}

\item \texttt{std::shuffle}将打乱\texttt{vector}中的内容，之后我们会再次对\texttt{vector}进行排序。前两个参数是容器的首尾迭代器，第三个参数是一个随机数生成器：

\begin{codebox}[cpp]
    shuffle(begin(v), end(v), g);
\end{codebox}

\item 现在\texttt{is\_sorted}函数将返回false，所以0将打印在屏幕上，\texttt{vector}的元素总量和具体数值都没有变，不过顺序发生了变化。我们会将函数的返回值再次打印在屏幕上：

\begin{codebox}[cpp]
    cout << is_sorted(begin(v), end(v)) << '\n';
    print(v);
\end{codebox}

\item 现在，在通过\texttt{std::sort}对\texttt{vector}进行排序。然后打印是否排序的结果：

\begin{codebox}[cpp]
    sort(begin(v), end(v));
    cout << is_sorted(begin(v), end(v)) << '\n';
    print(v);
\end{codebox}

\item 另一个比较有趣的函数是\texttt{std::partition}。有时候，并不需要对列表完全进行排序，只需要比它前面的某些值小就可以。所以，让使用\texttt{partition}将数值小于5的元素排到前面，并打印它们：

\begin{codebox}[cpp]
    shuffle(begin(v), end(v), g);
    partition(begin(v), end(v), [] (int i) { return i < 5; });
    print(v); 
\end{codebox}

\item 下一个与排序相关的函数是\texttt{std::partial\_sort}。我们可以使用这个函数对容器的内容进行排序，不过只是在某种程度上的排序。其会将\texttt{vector}中最小的N个数，放在容器的前半部分。其余的留在\texttt{vector}的后半部分，不进行排序：

\begin{codebox}[cpp]
    shuffle(begin(v), end(v), g);
    auto middle (next(begin(v), int(v.size()) / 2));
    partial_sort(begin(v), middle, end(v));
    print(v);
\end{codebox}

\item 当我们要对没做比较操作符的结构体进行比较，该怎么办呢？让我们来定义一个结构体，然后用这个结构体来实例化一个\texttt{vector}：

\begin{codebox}[cpp]
struct mystruct {
int a;
int b;
};
vector<mystruct> mv { {5, 100}, {1, 50}, {-123, 1000},
                   {3, 70}, {-10, 20} };
\end{codebox}

\item \texttt{std::sort}函数可以将比较函数作为第三个参数进行传入。让我们来使用它，并且传递一个比较函数。为了展示其实如何工作的，我们会对其第二个成员b进行比较。这样，我们将按\texttt{mystruct::b}的顺序进行排序，而非\texttt{mystruct::a}的顺序：

\begin{codebox}[cpp]
sort(begin(mv), end(mv),
[] (const mystruct &lhs, const mystruct &rhs) {
    return lhs.b < rhs.b;
});
\end{codebox}

\item 最后一步则是打印已经排序的\texttt{vector}：

\begin{codebox}[cpp]
    for (const auto &[a, b] : mv) {
        cout << "{" << a << ", " << b << "} ";
    }
    cout << '\n';
}
\end{codebox}

\item 编译运行程序。第一个1是由\texttt{std::is\_sorted}所返回的。之后将\texttt{vector}进行打乱后，\texttt{is\_sorted}就返回0。第三行是打乱后的\texttt{vector}。下一个1是使用sort之后进行打印的。然后，\texttt{vector}会被再次打乱，并且使用\texttt{std::partition}对部分元素进行排序。我们可以看到所有比5小的元素都在左边，比5大的都在右边。我们暂且将现在的顺序认为是乱序。倒数第二行展示了\texttt{std::partial\_sort}的结果。前半部分的内容进行了严格的排序，而后半部分则没有。最后一样，我们将打印\texttt{mystruct}实例的结果。其结果是严格根据第二个成员变量的值进行排序的：

\begin{codebox}[cpp]
$ ./sorting_containers
1
0
7, 1, 4, 6, 8, 9, 5, 2, 3, 10,
1
1, 2, 3, 4, 5, 6, 7, 8, 9, 10,
1, 2, 4, 3, 5, 7, 8, 10, 9, 6,
1, 2, 3, 4, 5, 9, 8, 10, 7, 6,
{-10, 20} {1, 50} {3, 70} {5, 100} {-123, 1000}
\end{codebox}

\end{enumerate}

\mySmallsection{How it works...}

这里我们使用了很多与排序算法相关的函数：

\begin{table}[htbp]\centering
\begin{tabularx}{\textwidth}{lX}
\toprule
算法函数 & 作用 \\
\midrule
\href{https://zh.cppreference.com/w/cpp/algorithm/sort}{std::sort} & 接受一定范围的元素，并对元素进行排序。 \\
\href{https://zh.cppreference.com/w/cpp/algorithm/is_sorted}{std::is\_sorted} & 接受一定范围的元素，并判断该范围的元素是否经过排序。 \\
\href{https://zh.cppreference.com/w/cpp/algorithm/random_shuffle}{std::shuffle} & 类似于反排序函数；其接受一定范围的元素，并打乱这些元素。 \\
\href{https://zh.cppreference.com/w/cpp/algorithm/partial_sort}{std::partial\_sort} & 接受一定范围的元素和另一个迭代器，前两个参数决定排序的范围，后两个参数决定不排序的范围。 \\
\href{https://zh.cppreference.com/w/cpp/algorithm/partition}{std::partition} & 能够接受谓词函数。所有元素都会在谓词函数返回true时，被移动到范围的前端。剩下的将放在范围的后方。 \\
\bottomrule
\end{tabularx}
\end{table}

对于没有实现比较操作符的对象来说，想要排序就需要提供一个自定义的比较函数。其签名为\texttt{bool function\_name(const T \&lhs, const T \&rhs)}，并且在执行过程中无副作用。

当然排序还有其他类似\texttt{std::stable\_sort}的函数，其能保证排序后元素的原始顺序，\texttt{std::stable\_partition}也有类似的功能。

\begin{myInfoBox}
Note: \texttt{std::sort}对于排序有不同的实现。根据所提供的迭代器参数，其实现分为选择排序、插入排序、合并排序，对于元素数量较少的容器可以完全进行优化。在使用者的角度，我们通常都不需要了解这些。
\end{myInfoBox}


% --- recipe 3: chapter5-3-chinese.md ---
\section{从容器中删除指定元素}

复制、转换和过滤是对一段数据常做的操作。本节，我们将重点放在过滤元素上。

将过滤出的元素从数据结构中移除，或是简单的移除其中一个，但对于不同数据结构来说，操作上就完全不一样了。在链表中(比如\texttt{std::list})，只要将对应节点的指针进行变动就好。不过，对于连续存储的结构体来说(比如\texttt{std::vector}，\texttt{std::array}，还有部分\texttt{std::deque})，删除相应的元素时，将会有其他元素来替代删除元素的位置。当一个元素槽空出来后，那么后面所有的元素都要进行移动，来将这个空槽填满。这个听起来都很麻烦，不过本节中我们只是想要从字符串中移除空格，这个功能没有太多的工作量。

当我们定义了一个结构体时，我们是不会考虑如何将其元素进行删除的。当需要做这件事的时候，我们才会注意到。STL中的\texttt{std::remove}和\texttt{std::remove\_if}函数可以给我们提供帮助。

\mySmallsection{How to do it...}

我们将通过不同的方式将\texttt{vector}中的元素进行删除：

\begin{enumerate}
\item 包含必要的头文件，并声明所使用的命名空间。

\begin{codebox}[cpp]
#include <iostream>
#include <vector>
#include <algorithm>
#include <iterator>
using namespace std;
\end{codebox}

\item 一个简单的打印辅助函数，用来打印\texttt{vector}中的内容：

\begin{codebox}[cpp]
void print(const vector<int> &v)
{
    copy(begin(v), end(v), ostream_iterator<int>{cout, ", "});
    cout << '\n';
}
\end{codebox}

\item 我们将使用简单的整数对\texttt{vector}进行实例化。然后，对\texttt{vector}进行打印，这样就能和后面的结果进行对比：

\begin{codebox}[cpp]
int main()
{
    vector<int> v {1, 2, 3, 4, 5, 6};
    print(v);
\end{codebox}

\item 现在，我们移除\texttt{vector}中所有的2。\texttt{std::remove}将2值移动到其他位置，这样这个值相当于消失了。因为\texttt{vector}长度在移除元素后变短了，\texttt{std::remove}将会返回一个迭代器，这个迭代器指向新的末尾处。旧的\texttt{end}迭代器所指向的地方，实际上就没有什么意义了，所以我们可以告诉\texttt{vector}将这个位置进行擦除。我们使用两行代码来完成这个任务：

\begin{codebox}[cpp]
{
    const auto new_end (remove(begin(v), end(v), 2));
    v.erase(new_end, end(v));
}
print(v);
\end{codebox}

\item 现在，我们来移除奇数。为了完成移除，我们需要实现一个谓词函数，这个函数用来告诉程序哪些值是奇数，然后结合\texttt{std::remove\_if}来使用。

\begin{codebox}[cpp]
{
    auto odd_number ([](int i) { return i % 2 != 0; });
    const auto new_end (
        remove_if(begin(v), end(v), odd_number));
    v.erase(new_end, end(v));
}
print(v);
\end{codebox}

\item 下一个尝试的算法是\texttt{std::replace}。我们使用这个函数将所有4替换成123。与\texttt{std::replace}函数对应，\texttt{std::replace\_if}也存在于STL中，同样可以接受谓词函数：

\begin{codebox}[cpp]
replace(begin(v), end(v), 4, 123);
print(v);
\end{codebox}

\item 让我们重新初始化\texttt{vector}，并为接下来的实验创建两个空的\texttt{vector}：

\begin{codebox}[cpp]
    v = {1, 2, 3, 4, 5, 6, 7, 8, 9, 10};
    vector<int> v2;
    vector<int> v3;
\end{codebox}

\item 然后，我们实现两个判读奇偶数的谓词函数：

\begin{codebox}[cpp]
    auto odd_number ([](int i) { return i % 2 != 0; });
    auto even_number ([](int i) { return i % 2 == 0; });
\end{codebox}

\item 接下来的两行做的事情完全相同。其将偶数拷贝到v2和v3中。第一行使用\texttt{std::remove\_copy\_if}函数，当相应数值不满足谓词条件时，函数会从源容器中拷贝到另一个容器中。第二行\texttt{std::copy\_if}则是拷贝满足谓词条件的元素。

\begin{codebox}[cpp]
remove_copy_if(begin(v), end(v),
    back_inserter(v2), odd_number);
copy_if(begin(v), end(v),
    back_inserter(v3), even_number); 
\end{codebox}

\item 然后，打印这两个\texttt{vector}，其内容应该是完全相同的。

\begin{codebox}[cpp]
    print(v2);
    print(v3);
}
\end{codebox}

\item 编译运行程序。第一行输出的是\texttt{vector}初始化的值。第二行是移除2之后的内容。接下来一行是移除所有奇数后的结果。第4行是将4替换为123的结果。最后两行则是v2和v3中的内容：

\begin{codebox}[cpp]
$ ./removing_items_from_containers
1, 2, 3, 4, 5, 6,
1, 3, 4, 5, 6,
4, 6,
123, 6,
2, 4, 6, 8, 10,
2, 4, 6, 8, 10,
\end{codebox}

\end{enumerate}

\mySmallsection{How it works...}

这里我们使用了很多与排序算法相关的函数：

\begin{table}[htbp]\centering
\begin{tabularx}{\textwidth}{lX}
\toprule
算法函数 & 作用 \\
\midrule
\href{https://zh.cppreference.com/w/cpp/algorithm/remove}{std::remove} & 接受一个容器范围和一个具体的值作为参数，并且移除对应的值。返回一个新的end迭代器，用于修改容器的范围。 \\
\href{https://zh.cppreference.com/w/cpp/algorithm/replace}{std::replace} & 接受一个容器范围和两个值作为参数，将使用第二个数值替换所有与第一个数值相同的值。 \\
\href{https://zh.cppreference.com/w/cpp/algorithm/remove_copy}{std::remove\_copy} & 接受一个容器范围，一个输出迭代器和一个值作为参数。并且将所有不满足条件的元素拷贝到输出迭代器的容器中。 \\
\href{https://zh.cppreference.com/w/cpp/algorithm/replace_copy}{std::replace\_copy} & 与\texttt{std::replace}功能类似，但与\texttt{std::remove\_copy}更类似些。源容器的范围并没有变化。 \\
\href{https://zh.cppreference.com/w/cpp/algorithm/copy}{std::copy\_if} & 与\texttt{std::copy}功能相同，可以多接受一个谓词函数作为是否进行拷贝的依据。 \\
\bottomrule
\end{tabularx}
\end{table}

\begin{myInfoBox}
Note: 表中没有if的算法函数，都有一个*\_if版本存在，其能接受谓词函数，通过谓词函数判断的结果来进行相应的操作。
\end{myInfoBox}


% --- recipe 4: chapter5-4-chinese.md ---
\section{改变容器内容}

如果说\texttt{std::copy}是STL中最简单的算法，那么\texttt{std::transform}就是第二简单的算法。和\texttt{copy}类似，其可将容器某一范围的元素放置到其他容器中，在这个过程中允许进行一些变换(变换函数会对输入值进行一定操作，然后再赋给目标容器)。此外，两个具有不同元素类型的容间也可以使用这个函数。这个函数超级简单，并且非常有用，这个函数会让标准组件具有更好的可移植性。

\mySmallsection{How to do it...}

本节，我们将使用\texttt{std::transform}在拷贝的同时，修改\texttt{vector}中的元素：

\begin{enumerate}
\item 包含必要的头文件，并且声明所使用的命名空间：

\begin{codebox}[cpp]
#include <iostream>
#include <vector>
#include <string>
#include <sstream>
#include <algorithm>
#include <iterator>
using namespace std;
\end{codebox}

\item \texttt{vector}由简单的整数组成：

\begin{codebox}[cpp]
int main()
{
    vector<int> v {1, 2, 3, 4, 5};
\end{codebox}

\item 为了打印元素，会将所有元拷贝素到\texttt{ostream\_iterator}适配器中。\texttt{transform}函数可以接受一个函数对象，其能在拷贝过程中对每个元素进行操作。这个例子中，我们将计算每个元素的平方值，所以代码将打印出平方数。因为直接进行了打印，所以平方数并没有进行保存：

\begin{codebox}[cpp]
transform(begin(v), end(v),
    ostream_iterator<int>{cout, ", "},
    [] (int i) { return i * i; });
cout << '\n';
\end{codebox}

\item 再来做另一个变换。例如，对于数字3来说，显示成\texttt{3\textasciicircum{}2 = 9}显然有更好的可读性。下面的辅助函数\texttt{int\_to\_string}函数对象就会使用\texttt{std::stringstream}对象进行打印操作：

\begin{codebox}[cpp]
auto int_to_string ([](int i) {
    stringstream ss;
    ss << i << "^2 = " << i * i;
    return ss.str();
});
\end{codebox}

\item 这样就可以将整型值放入字符串中。可以说我么将这个证书映射到字符串中。使用\texttt{transform}函数，使我们可以拷贝所有数值到一个字符串\texttt{vector}中：

\begin{codebox}[cpp]
vector<string> vs;
transform(begin(v), end(v), back_inserter(vs),
    int_to_string);
\end{codebox}

\item 在打印完成后，我们的例子就结束了：

\begin{codebox}[cpp]
    copy(begin(vs), end(vs),
        ostream_iterator<string>{cout, "\n"});
}
\end{codebox}

\item 编译并运行程序：

\begin{codebox}[cpp]
$ ./transforming_items_in_containers
1, 4, 9, 16, 25,
1^2 = 1
2^2 = 4
3^2 = 9
4^2 = 16
5^2 = 25
\end{codebox}

\end{enumerate}

\mySmallsection{How it works...}

\texttt{std::transform}函数工作原理和\texttt{std::copy}差不多，不过在拷贝的过程中会对源容器中的元素进行变换，这个变换函数由用户提供。


% --- recipe 5: chapter5-5-chinese.md ---
\section{在有序和无序的vector中查找元素}

通常，需要确定某种元素在某个容器范围内是否存在。如果存在，我们会对这个值进行修改，或者访问与其相关的值。

查找元素的目的是不同的。当想要让在一段已排序的元素中进行查找，可以使用二分查找法，这种方法要比线性查找快的多。如果没有排序，那么就只能进行线性遍历来查找对应的值。

传统的STL查找算法我们都可以使用，所以了解一下这些算法。本节将会使用两个不同的算法，线性查找算法\texttt{std::find}，二分查找算法\texttt{std::equal\_range}。

\mySmallsection{How to do it...}

本节，我们将对一个比较小的数据集进行线性和二分查找：

\begin{enumerate}
\item 包含必要的头文件，以及声明所使用的命名空间。

\begin{codebox}[cpp]
#include <iostream>
#include <vector>
#include <list>
#include <algorithm>
#include <string>
using namespace std;
\end{codebox}

\item 数据集会包含\texttt{city}结构体，只是存储的城市的名字和人口数量：

\begin{codebox}[cpp]
struct city {
    string name;
    unsigned population;
};  
\end{codebox}

\item 搜索算法需要将元素与目标对象进行对比，所以我们需要重载\texttt{city}结构体的\texttt{==}操作符：

\begin{codebox}[cpp]
bool operator==(const city &a, const city &b) {
    return a.name == b.name && a.population == b.population;
}
\end{codebox}

\item 我们也需要将\texttt{city}实例进行打印，所以我们对其输出操作符\texttt{<<}也进行了重载：

\begin{codebox}[cpp]
ostream& operator<<(ostream &os, const city &city) {
    return os << "{" << city.name << ", "
        << city.population << "}";
}
\end{codebox}

\item 查找函数通常会返回迭代器。当函数找到相应的元素时，会返回指向其的迭代器，否则就会返回容器的\texttt{end}迭代器。第二种情况下，我们就不能对该迭代器进行访问。因为要打印输出结果，所以需要实现一个函数，这个函数会返回另一个函数对象，并会将数据结构的\texttt{end}迭代器进行包装。当要对结果进行打印时，会与容器的\texttt{end}迭代器相比较，如果不是\texttt{end}，那么打印出查找到的值；如果是\texttt{end}，则仅打印\texttt{<end>}：

\begin{codebox}[cpp]
template <typename C>
static auto opt_print (const C &container)
{
    return [end_it (end(container))] (const auto &item) {
        if (item != end_it) {
            cout << *item << '\n';
        } else {
            cout << "<end>\n";
        }
    };
}
\end{codebox}

\item 我们使用德国的一些城市对\texttt{vector}进行实例化：

\begin{codebox}[cpp]
int main()
{
    const vector<city> c {
        {"Aachen", 246000},
        {"Berlin", 3502000},
        {"Braunschweig", 251000},
        {"Cologne", 1060000}
    };
\end{codebox}

\item 使用这个辅助函数构造一个城市打印函数，其会获取到城市\texttt{vector}容器的\texttt{end}迭代器\texttt{c}：

\begin{codebox}[cpp]
    auto print_city (opt_print(c));
\end{codebox}

\item 使用\texttt{std::find}在\texttt{vector}中找到相应的元素——科隆(Cologne)。因为可以直接获得这个元素，所以这个搜索看起来毫无意义。不过，在查找之前并不知道这个元素在\texttt{vector}中的位置，而\texttt{find}函数告诉我们这个元素的具体位置。我们也可以写一个循环，仅对城市名进行比较，而无需比较人口数量。不过，这是个不是很好的设计。下一步，我们将做另外一个实验：

\begin{codebox}[cpp]
{
    auto found_cologne (find(begin(c), end(c),
        city{"Cologne", 1060000}));
    print_city(found_cologne);
}
\end{codebox}

\item 当不需要知道对应城市的人口数量时，就不需要使用\texttt{==}操作符，只需要比较城市名称就好。\texttt{std::find\_if}函数可以接受一个函数对象作为谓词函数。这样，就能只使用城市名来查找“科隆”了：

\begin{codebox}[cpp]
{
    auto found_cologne (find_if(begin(c), end(c),
        [] (const auto &item) {
        return item.name == "Cologne";
        }));
    print_city(found_cologne);
}
\end{codebox}

\item 为了让搜索更加优雅，可以实现谓词构建器。\texttt{population\_higher\_than}函数对象能接受一个人口数量，并且返回人口数量比这个数量多的城市。在这个小数据集中找一下多于2百万人口的城市。例子中，只有柏林(Berlin)符合条件：

\begin{codebox}[cpp]
{
    auto population_more_than ([](unsigned i) {
        return [=] (const city &item) {
            return item.population > i;
        };
    });
    auto found_large (find_if(begin(c), end(c),
        population_more_than(2000000)));
    print_city(found_large);
}
\end{codebox}

\item 使用的查找函数，线性的遍历容器，查找的时间复杂度为O(n)。STL也有二分查找函数，其时间复杂度为O(log(n))。让我们生成一个新的数据集，其包含了一些整数，并构建了另一个\texttt{print}函数:

\begin{codebox}[cpp]
const vector<int> v {1, 2, 3, 4, 5, 6, 7, 8, 9, 10};
auto print_int (opt_print(v));
\end{codebox}

\item \texttt{std::binary\_search}函数会返回一个布尔值，这个布尔值会告诉你函数是否找到了相应的元素，但是不会将指向元素的迭代器返回。二分查找需要查找的列表是已排序的，否则二分查找将出错：

\begin{codebox}[cpp]
    bool contains_7 {binary_search(begin(v), end(v), 7)};
    cout << contains_7 << '\n';
\end{codebox}

\item 如果需要获得查找的元素，就需要使用其他STL函数。其中之一就是\texttt{std::equal\_range}。其不会返回对应元素的迭代器给我们，不过会返回一组迭代器。第一个迭代器是指向第一个不小于给定值的元素。第二个迭代器指向第一个大于给定值的元素。我们的范围为数字1到10，那么第一个迭代器将指向7，因为其是第一个不小于7的元素。第二个迭代器指向8，因为其实第一个大于7的元素：

\begin{codebox}[cpp]
    auto [lower_it, upper_it] (
        equal_range(begin(v), end(v), 7));
    print_int(lower_it);
    print_int(upper_it); 
\end{codebox}

\item 当需要其中一个迭代器，可以使用\texttt{std::lower\_bound}或\texttt{std::upper\_bound}。\texttt{lower\_bound}函数只会返回第一个迭代器，而\texttt{upper\_bound}则会返回第二个迭代器：

\begin{codebox}[cpp]
    print_int(lower_bound(begin(v), end(v), 7));
    print_int(upper_bound(begin(v), end(v), 7));
}
\end{codebox}

\item 编译并运行这个程序，我们看到如下输出：

\begin{codebox}[cpp]
$ ./finding_items
{Cologne, 1060000}
{Cologne, 1060000}
{Berlin, 3502000}
1
7
8
7
8
\end{codebox}

\end{enumerate}

\mySmallsection{How it works...}

本节使用的STL查找算法：

\begin{table}[htbp]\centering
\begin{tabularx}{\textwidth}{lX}
\toprule
算法函数 & 作用 \\
\midrule
\href{https://zh.cppreference.com/w/cpp/algorithm/find}{std::find} & 可将一个搜索范围和一个值作为参数。函数将返回找到的第一个值的迭代器。线性查找。 \\
\href{https://zh.cppreference.com/w/cpp/algorithm/find}{std::find\_if} & 与\texttt{std::find}原理类似，不过其使用谓词函数替换比较值。 \\
\href{https://zh.cppreference.com/w/cpp/algorithm/binary_search}{std::binary\_search} & 可将一个搜索范围和一个值作为参数。执行二分查找，当找到对应元素时，返回true；否则，返回false。 \\
\href{https://zh.cppreference.com/w/cpp/algorithm/lower_bound}{std::lower\_bound} & 可将一个查找返回和一个值作为参数，并且执行二分查找，返回第一个不小于给定值元素的迭代器。 \\
\href{https://zh.cppreference.com/w/cpp/algorithm/upper_bound}{std::upper\_bound} & 与\texttt{std::lower\_bound}类似，不过会返回第一个大于给定值元素的迭代器。 \\
\href{https://zh.cppreference.com/w/cpp/algorithm/equal_range}{std::equal\_range} & 可将一个搜索范围和一个值作为参数，并且返回一对迭代器。其第一个迭代器和\texttt{std::lower\_bound}返回结果一样，第二个迭代器和\texttt{std::upper\_bound}返回结果一样。 \\
\bottomrule
\end{tabularx}
\end{table}

所有这些函数，都能接受一个自定义的比较函数作为可选参数传入。这样就可以自定义的进行查找，就如我们在本章做的那样。

来看一下\texttt{std::equal\_range}是如何工作的。假设我们有一个\texttt{vector}，\texttt{ v = \{0, 1, 2, 3, 4, 5, 6, 7, 7, 7, 8\}}，并且调用\texttt{equal\_range(begin(v), end(v), 7);}，为了执行对7的二分查找。如\texttt{equal\_range}要返回一对上下限迭代器那样，这个结果将返回一段区域\texttt{\{7, 7, 7\}}，因为原始\texttt{vector}中有很多7，所以这个子队列中也有很多7。下图能说明其运行的原理：

\myGraphic{0.9}{images/chapter5/5-5-1.png}{}

首先，\texttt{equal\_range}会使用典型的二分查找，直到其找到那个不小于查找值的那个元素。而后，另一个迭代器也是用同样的方式找到。如同分开调用\texttt{lower\_bound}和\texttt{upper\_bound}一样。

为了获得一个二分查找函数，并返回其第一个适配条件的元素。我们可以按照如下的方式实现：

\begin{codebox}[cpp]
template <typename Iterator, typename T>
Iterator standard_binary_search(Iterator it, Iterator end_it, T value)
{
    const auto potential_match (lower_bound(it, end_it, value));
    if (potential_match != end_it && value == *potential_match) {
        return potential_match;
    }
    return end_it;
}
\end{codebox}

这个函数使用\texttt{std::lower\_bound}，为的就是找到第一个不大于\texttt{value}的元素。返回结果\texttt{potential\_match}，有三种情况：

\begin{itemize}
\item 没有值不小于\texttt{value}。这样，返回值和\texttt{end\_it}(\texttt{end}迭代器)一样。

\item 遇到的第一个不小于\texttt{value}的元素，同时也大于\texttt{value}。因此需要返回\texttt{end\_it}，表示没有找到相应的值。

\item \texttt{potential\_match}指向的元素与\texttt{value}相同。这个匹配没毛病。因此就返回相应的迭代器。

\end{itemize}

当类型T没有\texttt{==}操作符时，需要为二分查找提供一个\texttt{<}操作实现。然后，可以将比较重写为\texttt{!(value < *potential\_match) \&\& !(*potential\_match < value)}。如果它们不小于，也不大于，那么必定等于。

STL中因为缺少对多次命中的“定义”，所以并没有提供相应的函数来适配多次命中。

\begin{myInfoBox}
Note： 需要留意\texttt{std::map}和\texttt{std::set}等数据结构，它们有自己的\texttt{find}函数。它们携带的\texttt{find}函数要比通用的算法快很多，因为他们的实现与数据结构强耦合。
\end{myInfoBox}


% --- recipe 6: chapter5-6-chinese.md ---
\section{将vector中的值控制在特定数值范围内——std::clamp}

很多应用中，需要获得相应的数据。在对其进行绘制或进行其他处理前，会先对这些数据进行归一化，因为这些数据的差距很大。

通常可以使用\texttt{std::transform}通过传入一个谓词函数，对数据结构中的所有数据进行处理。不过，当不知道这些值有多大时或多小时，需要通过相应的函数找到数值的范围。

STL就包含这样的函数，比如\texttt{std::minmax\_element}和\texttt{std::clamp}。将这些函数与Lambda函数相结合，可以解决一些简单的任务。

\mySmallsection{How to do it...}

本节，将\texttt{vector}中的值使用两种不同的方式进行归一化，一种使用\texttt{std::minmax\_element}，另一种使用\texttt{std::clamp}：

\begin{enumerate}
\item 包含必要的头文件，并声明所使用的命名空间。

\begin{codebox}[cpp]
#include <iostream>
#include <vector>
#include <algorithm>
#include <iterator>
using namespace std;
\end{codebox}

\item 将实现一个获得最大值和最小值的函数。这里最大值和最小值会更新，以便我们进行处理。函数对象会获取最大最小值，并返回另一个函数对象，这个返回的函数对象会做一些实际的转换。为了简单起见，新的最小值为0，所以旧值不需要进行偏移，并且值的归一化都是相对于0。为了有更好的可读性，这里忽略了最大值和最小值可能是一个值的可能性，不过在实际程序中需要格外注意这点，否则就会遇到除零问题：

\begin{codebox}[cpp]
static auto norm (int min, int max, int new_max)
{
    const double diff (max - min);
    return [=] (int val) {
        return int((val - min) / diff * new_max);
    };
}
\end{codebox}

\item 另一个函数对象构造器成为\texttt{clampval}，其会返回一个函数对象用于捕获最小值和最大值，并调用\texttt{std::clamp}将值控制在一定范围内：

\begin{codebox}[cpp]
static auto clampval (int min, int max)
{
    return [=] (int val) -> int {
        return clamp(val, min, max);
    };
}
\end{codebox}

\item \texttt{vector}中需要归一化的值大小不一。这些数据可能是热度数据、海拔高度或股票金额：

\begin{codebox}[cpp]
int main()
{
    vector<int> v {0, 1000, 5, 250, 300, 800, 900, 321};
\end{codebox}

\item 为对这些值进行归一化，我们需要找到这个\texttt{vector}中的最大值和最小值。\texttt{std::minmax\_element}函数将帮助我们获得这两个值。其会返回一组迭代器来代表这两个值：

\begin{codebox}[cpp]
    const auto [min_it, max_it] (
        minmax_element(begin(v), end(v)));
\end{codebox}

\item 我们会将所有值从第一个\texttt{vector}拷贝到另一个中。让我们实例化第二个\texttt{vector}，并且让其接收第一个\texttt{vector}中的值：

\begin{codebox}[cpp]
vector<int> v_norm;
v_norm.reserve(v.size());
\end{codebox}

\item 使用\texttt{std::transform}从第一个\texttt{vector}拷贝到第二个\texttt{vector}。拷贝过程中，将会使用到归一化辅助函数。之前的最大值和最小值为0和1000。在归一化之后，为0和255：

\begin{codebox}[cpp]
transform(begin(v), end(v), back_inserter(v_norm),
    norm(*min_it, *max_it, 255));
\end{codebox}

\item 在实现另一个归一化策略之前，先将这个操作过后的结果进行打印：

\begin{codebox}[cpp]
    copy(begin(v_norm), end(v_norm),
        ostream_iterator<int>{cout, ", "});
    cout << '\n'; 
\end{codebox}

\item 对已经归一化的\texttt{vector}使用\texttt{clampval}，这时的最大值和最小值分别为255和0：

\begin{codebox}[cpp]
    transform(begin(v), end(v), begin(v_norm),
        clampval(0, 255));
\end{codebox}

\item 完成之后，打印所有元素：

\begin{codebox}[cpp]
    copy(begin(v_norm), end(v_norm),
        ostream_iterator<int>{cout, ", "});
    cout << '\n';
}   
\end{codebox}

\item 编译并运行程序。当前值的范围都在0到255之间，我们可以将其认为是RGB颜色的亮度值：

\begin{codebox}[cpp]
$ ./reducing_range_in_vector
0, 255, 1, 63, 76, 204, 229, 81,
0, 255, 5, 250, 255, 255, 255, 255,
\end{codebox}

\item 我们将对应的数据进行绘制，就得到了如下的图像。我们可以看到，使用最大最小值对原始数据进行变换，得到的数据时线性的。clamp曲线会损失一些信息。两种不同的结果在不同的情况下会很有用：

\myGraphic{0.9}{images/chapter5/5-6-1.png}{}

\end{enumerate}

\mySmallsection{How it works...}

除了\texttt{std::transform}，我们使用量两个算法：

\texttt{std::minmax\_element}能接受一对\texttt{begin}和\texttt{end}迭代器作为输入。其会对这个范围进行遍历，然后找到这个范围内的最大值和最小值。其返回值是一个组对，我们会在我们的缩放函数中使用这个组对。

\texttt{std::clamp}函数无法对一个范围进行可迭代操作。其接受三个值作为参数：一个给定值，一个最小值，一个最大值。这个函数的返回值则会将对应的值截断在最大值和最小值的范围内。我们也能使用\texttt{max(min\_val, min(max\_val, x))}来替代\texttt{ std::clamp(x, min\_val, max\_val) }。


% --- recipe 7: chapter5-7-chinese.md ---
\section{在字符串中定位模式并选择最佳实现——std::search}

在一个字符串中查找另一个字符串，与在一个范围内查找一个对象有些不同。首先，字符串是可迭代的对象。另一方面，从一个字符串中查询另一个字符串，就意味着就是在一个范围内查询另一个范围。所以在查找过程中，有多次的比较，所以我们需要其他算法参与。

\texttt{std::string}就包含\texttt{find}函数，其能实现我们想要的；不过，本节我们将使用\texttt{std::search}来完成这个任务。虽然，\texttt{std::search}在字符串中会大量的用到，不过很多种容器都能使用这个算法来完成查找任务。C++17之后，\texttt{std::search}添加了更多有趣的特性，并且其本身可使用简单地交换搜索算法。这些算法都优化过，并且免费提供给开发者使用。另外，我们可以实现自己的搜索算法，并且可以将我们实现的算法插入\texttt{std::search}中。

\mySmallsection{How to do it...}

我们将对字符串使用新\texttt{std::search}函数，并且尝试使用其不同的查找对象进行应用：

\begin{enumerate}
\item 首先，包含必要的头文件，和声明所要使用的命名空间。

\begin{codebox}[cpp]
#include <iostream>
#include <string>
#include <algorithm>
#include <iterator>
#include <functional>
using namespace std;
\end{codebox}

\item 我们将实现一个辅助函数，用于打印查找算法所范围的位置，从而输出子字符串。

\begin{codebox}[cpp]
template <typename Itr>
static void print(Itr it, size_t chars)
{
    copy_n(it, chars, ostream_iterator<char>{cout});
    cout << '\n';
}
\end{codebox}

\item 我们例子输入的一个勒庞风格的字符串，其中包含我们要查找的字符串。本例中，这个需要查找的字符串为"elitr":

\begin{codebox}[cpp]
int main()
{
    const string long_string {
        "Lorem ipsum dolor sit amet, consetetur"
        " sadipscing elitr, sed diam nonumy eirmod"};
    const string needle {"elitr"};
\end{codebox}

\item 旧\texttt{std::search}接口接受一组\texttt{begin}和\texttt{end}迭代器，用于确定子字符串的查找范围。这个接口会返回一个迭代器指向所查找到的子字符串。如果接口没有找到对应的字符串，其将返回该范围的\texttt{end}迭代器：

\begin{codebox}[cpp]
{
    auto match (search(begin(long_string), end(long_string),
                     begin(needle), end(needle)));
    print(match, 5);
}
\end{codebox}

\item C++17版本的\texttt{std::search}将会使用一组\texttt{begin/end}迭代器和一个所要查找的对象。\texttt{std::default\_searcher}能接受一组子字符串的\texttt{begin}和\texttt{end}迭代器，再在一个更大的字符串中，查找这个字符串：

\begin{codebox}[cpp]
{
    auto match (search(begin(long_string), end(long_string),
        default_searcher(begin(needle), end(needle))));
    print(match, 5);
}
\end{codebox}

\item 这种改变就很容易切换搜索算法。\texttt{std::boyer\_moore\_searcher}使用Boyer-Moore查找算法进行快速的查找：

\begin{codebox}[cpp]
{
    auto match (search(begin(long_string), end(long_string),
        boyer_moore_searcher(begin(needle), end(needle))));
    print(match, 5);
}   
\end{codebox}

\item C++17标准中，有三种不同的搜索器对象实现。其中还有一种是Boyer-Moore-Horspool查找算法实现：

\begin{codebox}[cpp]
    {
        auto match (search(begin(long_string), end(long_string),
            boyer_moore_horspool_searcher(begin(needle),
            end(needle))));
        print(match, 5);
    }
}
\end{codebox}

\item 我们编译并运行这个程序。我们可以看到相同的字符串输出：

\begin{codebox}[cpp]
$ ./pattern_search_string
elitr
elitr
elitr
elitr
\end{codebox}

\end{enumerate}

\mySmallsection{How it works...}

我们在\texttt{std::search}中使用了4种查找方式，得到了相同的结果。这几种方式适用于哪种情况呢？

让我们假设大字符串为\texttt{s}，要查找的部分为\texttt{p}。然后，调用\texttt{std::search(begin(s), end(s), begin(p), end(p));}和\texttt{std::search(begin(s), end(s), default\_searcher(begin(p), end(p));}做相同的事情。

其他搜索方式将会以更复杂的方式实现：

\begin{itemize}
\item \texttt{std::default\_searcher}：其会重定向到\texttt{std::search}的实现。

\item \texttt{std::boyer\_moore\_searcher}：使用Boyer-Moore查找算法。

\item \texttt{std::boyer\_moore\_horspool\_searcher}：使用Boyer-Moore-Horspool查找算法。

\end{itemize}

为什么会有这些特殊的算法呢？Boyer-Moore算法起源于一个特殊想法——查找部分与原始字符串进行比较，其始于查找字符串的尾部，并且从右到左查找。如果查找的字符多个位置不匹配，并且对应部分没有出现，那么就需要在整个字符串进行平移，然后在进行查找。下图可能会看的更加明白一些。先来看一下第一步发生了什么：因为算法已知所要匹配字符串的长度，所以需要对相应位置上的字符进行比较，然后在平移到下一个长度点进行比较。在图中，这发生在第二步。这样Boyer-Moore算法就能避免对不必要的字符进行比较。

\myGraphic{0.9}{images/chapter5/5-7-1.png}{}

当然，在我们没有提供新的比较查找算法时，Boyer-Moore为默认的查找算法。其要比原先默认的算法快很多，不过其需要快速查找的数据结果进行支持，以判断搜索字符是否存在于查找块中，以及以多少为偏移进行定位。编译器将选择不同复杂度的算法实现，这取决于其所使用到的数据类型(对复杂类型的哈希映射和类型的原始查找表进行更改)。最后，默认的查找算法在查询不是很长的字符串也是非常的快。如果需要查找算法提高性能，那么Boyer-Moore将会是个不错的选择。

Boyer-Moore-Horspool为简化版的Boyer-Moore算法。其丢弃了“坏字符”规则，当对应字符串没有找到时，将会对整个查找块进行偏移。需要权衡的是，这个算法要比Boyer-Moore算法慢，但是其不需要对那么多特殊的数据结构进行操作。

\begin{myInfoBox}
Note： 不要试图尝试比较哪种算法在哪种情况下更快。你可以使用自己实际的例子进行测试，并且基于你得到的结果进行讨论。
\end{myInfoBox}


% --- recipe 8: chapter5-8-chinese.md ---
\section{对大vector进行采样}

有时我们需要处理非常庞大的数据量，不可能在短时间内处理完这些数据。这样的话，数据可能就需要采样来减少要处理的数据量，从而加速整个处理过程。另一些情况下，不减少数据量也能加快程序处理的速度，不过这需要对一些数据进行存储或变换。

采样最原始的方式是每隔N个数据点，采样一次。在大多数情况下这样做没有问题，但是在信号处理中，其会引发一种称为\textbf{混淆}的数学情况。当减少两个随机采样点的距离时，这种现象会减弱。我们看一下下面的图，这张图就很能说明问题——当原始信号为一个sin波时，图例为三角的曲线就表示对这个曲线进行每隔100个点的取样。

\myGraphic{0.9}{images/chapter5/5-8-1.png}{}

不幸的是，其采样得到的值都是同一个Y值！连接起来就是与X轴平行的一条线。平方点采样，其每隔\texttt{100+random(-15, +15)}个值进行采样。不过，这样连接起来的曲线看起来和原始的曲线还是相差很远，所以在这个例子中就不能以固定的步长进行采样。

\texttt{std::sample}函数不会添加随机值来改变采样的步长，而是采用完全随机的点进行采样。所以其工作方式与上图所显示的大为不同。

\mySmallsection{How to do it...}

我们将对一个具有随机值的大vector进行采样。随机数据符合正态分布。采样结果也要符合正态分布，来让我们看下代码：

\begin{enumerate}
\item 首先包含必要的头文件，以及声明所使用的命名空间。

\begin{codebox}[cpp]
#include <iostream>
#include <vector>
#include <random>
#include <algorithm>
#include <iterator>
#include <map>
#include <iomanip>
using namespace std; 
\end{codebox}

\item 使用常数直接对变量进行初始化。第一个值代表了\texttt{vector}的的长度，第二个数代表了采样的步长：

\begin{codebox}[cpp]
int main()
{
    const size_t data_points {100000};
    const size_t sample_points {100};
\end{codebox}

\item 我们要使用符合正态分布的随机值生成器来将\texttt{vector}填满。这里先来确定正太分布的平均值和标准差：

\begin{codebox}[cpp]
    const int mean {10};
    const size_t dev {3};
\end{codebox}

\item 现在，我们来设置随机数生成器。首先，我们实例化一个随机设备，然后给定一个随机种子，对生成器进行初始化。然后，就可以得到对应分布的随机生成器：

\begin{codebox}[cpp]
    random_device rd;
    mt19937 gen {rd()};
    normal_distribution<> d {mean, dev};
\end{codebox}

\item 对\texttt{vector}进行初始化，并用随机值将\texttt{vector}进行填充。这里会使用到\texttt{std::generate\_n}算法，其会将随机值，通过\texttt{back\_inserter}迭代器插入\texttt{vector}中。生成函数对象包装成了\texttt{d(gen)}表达式，其能生成符合分布的随机值：

\begin{codebox}[cpp]
    vector<int> v;
    v.reserve(data_points);
    generate_n(back_inserter(v), data_points,
        [&] { return d(gen); });
\end{codebox}

\item 我们再实例化另一个\texttt{vector}，其来放采样过后的数值：

\begin{codebox}[cpp]
    vector<int> samples;
    v.reserve(sample_points);
\end{codebox}

\item \texttt{std::sample}算法与\texttt{std::copy}的原理类似，不过其需要两个额外的参数：采样数量和随机值生成对象。前者确定输入范围，后者去确定采样点：

\begin{codebox}[cpp]
    sample(begin(v), end(v), back_inserter(samples),
        sample_points, mt19937{random_device{}()});
\end{codebox}

\item 这样就完成了采样。代码的最后展示一下我们的采样结果。输入数据符合正态分布，如果采样算法可行，那么其采样的结果也要符合正态分布。为了展示采样后的值是否符合正态分布，我们将数值的直方图进行打印：

\begin{codebox}[cpp]
    map<int, size_t> hist;
    for (int i : samples) { ++hist[i]; }
\end{codebox}

\item 最后，我们使用循环打印出直方图：

\begin{codebox}[cpp]
    for (const auto &[value, count] : hist) {
        cout << setw(2) << value << " "
            << string(count, '*') << '\n';
    }
}
\end{codebox}

\item 编译并运行程序，我们将看到采样后的结果，其也符合正态分布：\myGraphic{0.9}{images/chapter5/5-8-2.png}{}

\end{enumerate}

\mySmallsection{How it works...}

\texttt{std::sample}算法是C++17添加的。其函数签名如下：

\begin{codebox}[cpp]
template<class InIterator, class OutIterator,
class Distance, class UniformRandomBitGenerator>
OutIterator sample(InIterator first, InIterator last,
                 SampleIterator out, Distance n,
                 UniformRandomBitGenerator&& g);
\end{codebox}

其输入范围有first和last迭代器确定，\texttt{out}迭代器作为采样输出。这些迭代器对于该函数来说和\texttt{std::copy}类似，元素从一个容器拷贝到另一个。\texttt{std::sample}算法只会拷贝输入中的一部分，因为采样结果只有n个元素。其在内部使用均匀分布，所以能以相同的概率选择输入范围中的每个数据点。


% --- recipe 9: chapter5-9-chinese.md ---
\section{生成输入序列的序列}

当测试代码需要处理参数顺序不重要的输入序列时，有必要测试它是否对所有可能的输入产生相同的输出。当你自己实现了一个排序算法时，就要写这样的测试代码来确定自己的实现是否正确。

\texttt{std::next\_permutation}在任何时候都能帮我们将序列进行打乱。我们在可修改的范围中可以调用它，其会将以字典序进行置换。

\mySmallsection{How to do it...}

本节，我们将从标准输入中读取多个字符串，然后使用\texttt{std::next\_permutation}生成已排序的序列，并且打印这个序列：

\begin{enumerate}
\item 首先，包含必要的头文件，并声明所使用的命名空间。

\begin{codebox}[cpp]
#include <iostream>
#include <vector>
#include <string>
#include <iterator>
#include <algorithm>
using namespace std; 
\end{codebox}

\item 使用标准数组对\texttt{vector}进行初始化，接下来对\texttt{vector}进行排序：

\begin{codebox}[cpp]
int main()
{
    vector<string> v {istream_iterator<string>{cin}, {}};
    sort(begin(v), end(v));
\end{codebox}

\item 现在来打印\texttt{vector}中的内容。随后，调用\texttt{std::next\_permutation}，其会打乱已经排序的\texttt{vector}，再对其进行打印。直到\texttt{next\_permutation}返回false时，代表\texttt{next\_permutation}完成了其操作，循环结束：

\begin{codebox}[cpp]
    do {
        copy(begin(v), end(v),
            ostream_iterator<string>{cout, ", "});
        cout << '\n';
    } while (next_permutation(begin(v), end(v)));
}
\end{codebox}

\item 编译运行这个程序，会有如下的打印：

\begin{codebox}[cpp]
$ echo "a b c" | ./input_permutations
a, b, c,
a, c, b,
b, a, c,
b, c, a,
c, a, b,
c, b, a,
\end{codebox}

\end{enumerate}

\mySmallsection{How it works...}

\texttt{std::next\_permutation}算法使用起来有点奇怪。因为这个函数接受一组开始/结束迭代器，当其找到下一个置换时返回true；否则，返回false。不过“下一个置换”又是什么意思呢？

当\texttt{std::next\_permutation}算法找到元素中的下一个字典序时，其会以如下方式工作：

\begin{enumerate}
\item 通过\texttt{v[i - 1] < v[i]}的方式找到最大索引i。如果这个最大索引不存在，那么返回false。

\item 再找到最大所以j，这里j需要大于等于i，并且\texttt{v[j] > v[i - 1]}。

\item 将位于索引位置j和i - 1上的值进行交换。

\item 将从i到范围末尾的元素进行反向。

\item 返回true。

\end{enumerate}

每次单独的置换顺序，都会在同一个序列中呈现。为了看到所有置换的可能，我们先对数组进行了排序。如果我们输入“c b a”到算法中，算法会立即终止，因为每个元素都以反字典序排列。


% --- recipe 10: chapter5-10-chinese.md ---
\section{实现字典合并工具}

假设我们有一个已经排序的列表，有人有另一个已排序的列表，我们想要将这两个列表进行共享。那么最好的方式就是将这两个列表合并起来。我们需要合并后的列表也是有序的，这样我们查找元素就会十分方便。

为了将两个已排序列表中的元素进行合并，我们本能的会想需要创建一个新的列表来放置这两个列表中的元素。对于要加入的元素，我们需要将队列中的元素进行对比，然后找到最小的那个元素将其放到列表的最前面。不过，这样输出队列的顺序会被打乱。下面的图就能很好的说明这个问题：

\myGraphic{0.9}{images/chapter5/5-10-1.png}{}

\texttt{std::merge}算法就可以直接来帮助我们做这个事情，这样我们就无需过多的参与。本节我们将展示如何使用这个算法。

\mySmallsection{How to do it...}

我们将创建一个简单的字典，其为英语单词和德语单词一对一的翻译，之后将其存储在\texttt{std::deque}数据结构中。程序将标注输入中获取这个字典，并且打印合并之后的字典。

\begin{enumerate}
\item 包含必要的头文件，并声明所使用的命名空间。

\begin{codebox}[cpp]
#include <iostream>
#include <algorithm>
#include <iterator>
#include <deque>
#include <tuple>
#include <string>
#include <fstream>
using namespace std; 
\end{codebox}

\item 字典是一对字符串，两两对应：

\begin{codebox}[cpp]
using dict_entry = pair<string, string>;
\end{codebox}

\item 我们将在屏幕上打印这个组对，并且要从用户输入中读取这个组对，所以我们必须要重载\texttt{>>}和\texttt{<<}操作符：

\begin{codebox}[cpp]
namespace std {
ostream& operator<<(ostream &os, const dict_entry p)
{
    return os << p.first << " " << p.second;
}
istream& operator>>(istream &is, dict_entry &p)
{
    return is >> p.first >> p.second;
}
}
\end{codebox}

\item 这里需要创建一个辅助函数，其能接受任何流对象作为输入，帮助我们构建字典。其会构建一个\texttt{std::deque}来存放一对一的字符串对，并且其会读取标准输入中的所有字符。并在返回字典前，对字典进行排序：

\begin{codebox}[cpp]
template <typename IS>
deque<dict_entry> from_instream(IS &&is)
{
    deque<dict_entry> d {istream_iterator<dict_entry>{is}, {}};
    sort(begin(d), end(d));
    return d;
}
\end{codebox}

\item 这里使用不同的输入流，创建两个不同的字典。其中一个是从\texttt{dict.txt}文件中读取出的字符，我们先假设这个文件存在。其每一行为一个组对，另一个流就是标准输入：

\begin{codebox}[cpp]
int main()
{
    const auto dict1 (from_instream(ifstream{"dict.txt"}));
    const auto dict2 (from_instream(cin));
\end{codebox}

\item 作为辅助函数\texttt{from\_instream}将返回给我们一个已经排过序的字典，这样我们就可以将两个字典直接放入\texttt{std::merge}算法中。其能通过给定两个的\texttt{begin}和\texttt{end}迭代器组确定输入的范围，并在最后给定输出。这里的输出将会打印在用户的屏幕上：

\begin{codebox}[cpp]
    merge(begin(dict1), end(dict1),
        begin(dict2), end(dict2),
        ostream_iterator<dict_entry>{cout, "\n"});
}
\end{codebox}

\item 可以编译这个程序，不过在运行之前，我们需要创建\texttt{dict.txt}文件，并且写入如下内容：

\begin{codebox}[cpp]
car auto
cellphone handy
house haus
\end{codebox}

\item 现在我们运行程序了，输入一些英文单词，将其翻译为德文。这时的输出仍旧是一个排序后的字典，其可以将输入的所有单词进行翻译。

\begin{codebox}[cpp]
$ echo "table tisch fish fisch dog hund" | ./dictionary_merge
car auto
cellphone handy
dog hund
fish fisch
house haus
table tisch
\end{codebox}

\end{enumerate}

\mySmallsection{How it works...}

\texttt{std::meger}算法接受两对\texttt{begin/end}迭代器，这两对迭代器确定了输入范围。这两对迭代器所提供的输入范围也必须是已排序的。第五个参数就是输出容器的迭代器，其接受两段范围合并的元素。

其有一个变体\texttt{std::inplace\_merge}。两个算法几乎一样，不过这个变体只需要一对迭代器，并且没有输出，和其名字一样，其会直接在输入范围上进行操作。比如对\texttt{\{A, C, B, D\}}这个序列来说，可以将第一个子序列定义为\texttt{\{A, C\}}，第二个子序列定义为\texttt{\{B, D\}}。使用\texttt{std::inplace\_merge}算法将两个序列进行合并，其结果为\texttt{\{A, B, C, D\}}。


